% ==============================================================================
%  Paquetes y configuración general
% ==============================================================================

% --- Idioma y fuentes ----------------------------------------------------------
\usepackage[T1]{fontenc}
\usepackage[utf8]{inputenc}
\usepackage{lmodern}
\usepackage[spanish,es-nodecimaldot,es-tabla,es-noshorthands]{babel}
\usepackage{microtype}

% --- Matemáticas ---------------------------------------------------------------
\usepackage{amsmath, amssymb, amsthm, mathtools}
\usepackage{bm}

% --- Página --------------------------------------------------------------------
\usepackage[a4paper, top=2.5cm, bottom=2.5cm, left=2.7cm, right=2.7cm,
            headheight=14pt]{geometry}
\usepackage{setspace}
\setstretch{1.15}
\setlength{\parindent}{0pt}
\setlength{\parskip}{0.5em}

% --- Tablas, listas, figuras ---------------------------------------------------
\usepackage{booktabs}
\usepackage{array}
\usepackage{multirow}
\usepackage{enumitem}
\setlist{itemsep=0.2em, topsep=0.3em}
\usepackage{graphicx}
\graphicspath{{figuras/}}
\usepackage{float}
\usepackage{caption}
\usepackage{subcaption}

% --- Gráficos ------------------------------------------------------------------
\usepackage{tikz}
\usetikzlibrary{arrows.meta, positioning, shapes, calc, decorations.pathreplacing, babel}
\usepackage{pgfplots}
\pgfplotsset{compat=1.18}
\usepackage{pgfplotstable}

% --- Colores y cajas -----------------------------------------------------------
\usepackage[dvipsnames]{xcolor}
\usepackage[most]{tcolorbox}

% --- Encabezados y títulos ----------------------------------------------------
\usepackage{titlesec}
\usepackage{fancyhdr}

% --- Enlaces (cargar casi al final) ---------------------------------------------
\definecolor{enlace}{RGB}{50,75,120}
\usepackage[colorlinks=true, linkcolor=enlace, urlcolor=enlace,
            citecolor=enlace, bookmarksnumbered=true]{hyperref}
\usepackage{bookmark}
\usepackage[spanish, nameinlink, noabbrev]{cleveref}

% --- Estilo de capítulos (un capítulo = un tema) -------------------------------
\titleformat{\chapter}[display]
  {\normalfont\bfseries}
  {\LARGE Tema \thechapter}
  {0.4em}
  {\Huge}
  [{\vspace{0.3em}\titlerule[1pt]}]
\titlespacing*{\chapter}{0pt}{-20pt}{25pt}

\titleformat{\section}
  {\normalfont\Large\bfseries}{\thesection}{0.8em}{}
\titleformat{\subsection}
  {\normalfont\large\bfseries}{\thesubsection}{0.8em}{}

% --- Encabezado / pie ----------------------------------------------------------
\pagestyle{fancy}
\fancyhf{}
\renewcommand{\chaptermark}[1]{\markboth{Tema \thechapter. #1}{}\gdef\tituloTemaActual{#1}}
\fancyhead[L]{\small\nouppercase{\leftmark}}
\fancyhead[R]{\small Regresión y ANOVA}
\fancyfoot[C]{\thepage}
\renewcommand{\headrulewidth}{0.4pt}
\fancypagestyle{plain}{\fancyhf{}\fancyfoot[C]{\thepage}\renewcommand{\headrulewidth}{0pt}}

% --- Estilos de gráficos comunes -----------------------------------------------
\pgfplotsset{
  mini/.style={width=4.4cm, height=4.2cm, ticks=none, axis lines=box,
               xlabel={$X$}, ylabel={$Y$}, xlabel style={yshift=2pt},
               ylabel style={rotate=-90, yshift=-2pt}, every axis plot/.append style={thin}},
  puntos/.style={only marks, mark=*, mark size=0.9pt, black},
  recta/.style={NavyBlue, thick, no marks},
  % gráficos de residuos (Tema 4): como mini, con r_i frente a los valores ajustados
  residuos/.style={mini, xlabel={$\est{y}_i$}, ylabel={$r_i$}, xmin=0, xmax=10},
  cero/.style={gray, dashed, no marks, domain=0:10},
}
